% Hafta 2 — Bir açığın yaşam çizgisinde kim ne yapar?
% Derleme: py -3.12 tools/latex/derle.py docs/week-2/assets/h02-17-acik-yasam-cizgisi.tex
\documentclass[tikz,border=8pt]{standalone}
\usepackage{cen429-cizim}
\begin{document}
\begin{tikzpicture}[node distance=8mm and 8mm]
  \node[baslik] (b) at (0,0) {Bir açığın yaşam çizgisinde kim ne yapar?};
  \node[altbaslik] at (b.south west) {en tehlikeli aralık: yama çıktı ama kurulmadı};

  \node[morkutu=38mm, minimum height=18mm, below=13mm of b.south west, anchor=north west, xshift=8mm] (s1)
    {\kb{Keşif}\\araştırmacı\\bulur};
  \node[kutu=38mm, minimum height=18mm, right=of s1] (s2) {\kb{Özel ifşa}\\üreticiye\\bildirilir};
  \node[kutu=38mm, minimum height=18mm, right=of s2] (s3) {\kb{Yama geliştirme}\\üretici\\düzeltir};
  \draw[ok, draw=soluk] (s1) -- (s2);
  \draw[ok, draw=soluk] (s2) -- (s3);

  \node[iyikutu=38mm, minimum height=18mm, below=16mm of s1] (s4) {\kb{Yama yayını}\\CVE + bülten};
  \node[uyarikutu=38mm, minimum height=18mm, right=of s4] (s5) {\kb{Kamuya açılma}\\ayrıntı yayımlanır};
  \node[kotukutu=38mm, minimum height=18mm, right=of s5] (s6) {\kb{Kitlesel sömürü}\\yamasızlar vurulur};
  \draw[ok, draw=soluk] (s4) -- (s5);
  \draw[kotuok] (s5) -- (s6);
  \draw[ok, draw=soluk] (s3.south) to[out=-90, in=90] (s4.north);

  \node[kotudip, text width=180mm, below=10mm of s4.south -| s5, anchor=north]
    {WannaCry'ın dersi tam burada: yama saldırıdan haftalar önce vardı.};
\end{tikzpicture}
\end{document}
